\documentclass[11pt,a4paper]{article}

\usepackage[utf8]{inputenc}
\usepackage[T1]{fontenc}
\usepackage[spanish,es-tabla,es-noshorthands]{babel}
\usepackage{lmodern}
\usepackage{microtype}
\usepackage[a4paper,margin=2.5cm,headheight=14.5pt]{geometry}
\usepackage{setspace}\onehalfspacing
\usepackage{parskip}

\usepackage[table]{xcolor}
\definecolor{uteqverde}{RGB}{0,102,51}
\definecolor{uteqdorado}{RGB}{204,153,0}
\definecolor{grisfila}{RGB}{245,245,245}
\definecolor{goldclaro}{RGB}{252,245,220}
\definecolor{verdeclaro}{RGB}{235,247,238}
\definecolor{textogris}{RGB}{60,60,60}
\definecolor{nivelaus}{RGB}{176,42,42}
\definecolor{codebg}{RGB}{248,248,248}

\usepackage{amsmath,amssymb}
\usepackage{graphicx}
\usepackage{tikz}
\usepackage{fancyhdr}
\usepackage[most]{tcolorbox}

\newtcolorbox{cajadefinicion}[1]{
    colback=uteqverde!5,
    colframe=uteqverde,
    fonttitle=\bfseries,
    coltitle=white,
    title={#1},
    boxrule=0.6pt,
    arc=3pt,
    breakable
}

\newtcolorbox{cajaentregable}[1]{
    colback=uteqdorado!8,
    colframe=uteqdorado!80!black,
    fonttitle=\bfseries,
    coltitle=white,
    title={#1},
    boxrule=0.6pt,
    arc=3pt,
    breakable
}

\newtcolorbox{cajacritica}[1]{
    colback=red!5,
    colframe=red!70!black,
    fonttitle=\bfseries,
    coltitle=white,
    title={#1},
    boxrule=0.6pt,
    arc=3pt,
    breakable
}

\newtcolorbox{cajarepo}[1]{
    colback=blue!4,
    colframe=blue!60!black,
    fonttitle=\bfseries,
    coltitle=white,
    title={#1},
    boxrule=0.6pt,
    arc=3pt,
    breakable
}

\usepackage{listings}
\lstset{
    basicstyle=\ttfamily\footnotesize,
    backgroundcolor=\color{codebg},
    frame=single,
    breaklines=true,
    columns=fullflexible,
    keepspaces=true
}

\usepackage{array}
\usepackage{booktabs}
\usepackage{longtable}
\usepackage{tabularx}
\usepackage{multirow}

\usepackage[hidelinks,bookmarks=true,bookmarksnumbered=true]{hyperref}
\usepackage{enumitem}\setlist{itemsep=2pt,topsep=2pt}

\newcolumntype{C}[1]{>{\centering\arraybackslash}p{#1}}
\newcolumntype{L}[1]{>{\raggedright\arraybackslash}p{#1}}

\pagestyle{fancy}
\fancyhf{}
\lhead{\footnotesize\textcolor{uteqverde}{Aplicaciones Web --- Quinto nivel --- PPA 2026-2027}}
\rhead{\footnotesize\textcolor{uteqverde}{PFC --- Guia de la Tercera Entrega}}
\rfoot{\thepage}
\lfoot{\footnotesize UTEQ --- FCI --- Software}
\renewcommand{\headrulewidth}{0.4pt}

\hypersetup{
    pdftitle={Guia de la Tercera Entrega del PFC --- Aplicaciones Web},
    pdfauthor={Dr. Gleiston Guerrero Ulloa},
    pdfsubject={Tercera Entrega del Proyecto Fin de Curso},
    pdfkeywords={PFC, Reproducibilidad, FAIR, evidencia empirica, Aplicaciones Web, UTEQ}
}

\begin{document}

% =====================================================================
%  PORTADA
% =====================================================================
\begin{center}
    \vspace*{1.5cm}
    {\Large\bfseries\textcolor{uteqverde}{UNIVERSIDAD TECNICA ESTATAL DE QUEVEDO}}\\[4pt]
    {\large Facultad de Ciencias de la Computacion y Diseno Digital}\\[2pt]
    {\large Carrera de Ingenieria de Software}\\[24pt]

    {\LARGE\bfseries Guia de la Tercera Entrega del}\\[6pt]
    {\LARGE\bfseries Proyecto Fin de Curso (PFC)}\\[18pt]

    {\large Aplicaciones Web --- Quinto nivel}\\[2pt]
    {\large Periodo Academico Presencial 2026-2027}\\[24pt]

    \begin{tcolorbox}[colback=uteqverde!8,colframe=uteqverde,arc=3pt,
                      width=0.85\textwidth,boxrule=0.8pt]
        \centering
        \textbf{Docente:} Dr. Gleiston Cicer\'on Guerrero Ulloa, Ph.D.\\[2pt]
        \texttt{gguerrero@uteq.edu.ec}\\[6pt]
        \textbf{Etiqueta objetivo de esta entrega:} \texttt{v0.9.0-rc}\\[2pt]
        \textbf{Ubicacion temporal:} entre la Entrega 2 (\texttt{v0.7.0}, semana 15) y\\ la Entrega Final (\texttt{v1.0.0}, semana 16)
    \end{tcolorbox}

    \vfill

    {\small Quevedo --- Los Rios --- Ecuador}\\[2pt]
    {\small 2026}
\end{center}

\newpage

% =====================================================================
%  1. NATURALEZA
% =====================================================================
\section{Naturaleza y ubicacion temporal de la Tercera Entrega}

La Tercera Entrega del PFC ocupa el hito intermedio entre la Entrega 2 (\texttt{v0.7.0}, semana 15) y la Entrega Final (\texttt{v1.0.0}, semana 16). Su proposito no es cerrar nuevas funcionalidades sino consolidar y verificar, con evidencia empirica reproducible, todo lo entregado hasta la Entrega 2, elevando el proyecto al estado \emph{release candidate} (\texttt{v0.9.0-rc}) exigido por las normas de reproducibilidad de las revistas de alto impacto.

Mientras que la Entrega 2 verifica lo que el sistema \emph{hace} (API REST autenticada, CRUD Angular, cache Redis, ProblemDetails y pruebas unitarias basicas), la Tercera Entrega verifica \emph{como} lo hace, \emph{con que garantia} y \emph{con que grado de reproducibilidad por un tercero sin conocimiento previo del proyecto}. Esta separacion es lo que las normas contemporaneas de gestion de datos de investigacion identifican como el paso de un producto software a un artefacto de investigacion reutilizable~\cite{wilkinson2016fair,acm2020artifact}.

En consecuencia, la Tercera Entrega solicita al equipo tres transformaciones sobre el estado \texttt{v0.7.0} que ya posee.

\begin{enumerate}
    \item \textbf{Consolidacion de la reproducibilidad automatica}: el proyecto debe poder ser reconstruido de forma identica por un tercero desde una clonacion limpia, con contenedores pinados por \emph{digest} \texttt{sha256}, semillas deterministas fijadas y \emph{scripts} de un solo comando para levantar, ejecutar pruebas, generar metricas y apagar el sistema.
    \item \textbf{Consolidacion de la evidencia empirica cuantitativa}: cada afirmacion no trivial del sistema (rendimiento, seguridad, usabilidad, fiabilidad, cobertura) debe estar respaldada por un archivo de mediciones crudo, un \emph{script} de analisis y un reporte agregado con estadistica descriptiva e inferencial minima, todo bajo control de version.
    \item \textbf{Consolidacion de la trazabilidad}: cada requisito funcional y no funcional debe estar trazado hasta su codigo, su prueba automatizada y su evidencia empirica de cumplimiento, siguiendo la disciplina de \emph{decision records} propuesta por Nygard y adoptada por la ingenieria de arquitectura contemporanea~\cite{nygard2011adr,brown2023c4,bass2021software,richards2020fundamentals}.
\end{enumerate}

La lista siguiente resume la posicion de la Tercera Entrega en la linea temporal del PFC.

\begin{cajadefinicion}{Linea temporal de las entregas del PFC}
    \begin{itemize}
        \item \textbf{Entrega 1} (semana 5): planteamiento del problema, seleccion de la pila, esqueleto ejecutable con Docker, tag \texttt{v0.3.0}.
        \item \textbf{Entrega 2} (semana 15): arquitectura C4 documentada, API REST + JWT + Swagger, Angular con CRUD, ProblemDetails RFC 7807, cache Redis, cinco pruebas unitarias JUnit 5 + MockMvc, tag \texttt{v0.7.0}~\cite{brown2023c4,rfc7807,jones2015jwt}.
        \item \textbf{Entrega 3} (esta guia): consolidacion de reproducibilidad, evidencia empirica cuantitativa y trazabilidad, tag \texttt{v0.9.0-rc}.
        \item \textbf{Entrega Final} (semana 16): version estable con cobertura JaCoCo $\geq 70$\,\%, k6 y Lighthouse dentro de umbrales, auditoria de los seis controles OWASP minimos y sistema listo para archivo permanente, tag \texttt{v1.0.0}~\cite{owasp2021top10,ortega2022securecoding}.
    \end{itemize}
\end{cajadefinicion}

% =====================================================================
%  2. REQUISITOS
% =====================================================================
\section{Requisitos ineludibles de la Tercera Entrega}

Los requisitos que siguen consolidan lo entregado en la Entrega 2 y anaden las capas de evidencia empirica y de gestion de artefactos que las revistas de alto impacto exigen para aceptar un proyecto como reutilizable y reproducible por un tercero sin comunicacion con los autores originales~\cite{wilkinson2016fair,acm2020artifact,smith2016software}. La lectura del apartado debe hacerse en el orden en que aparece: cada bloque supone que el anterior ya esta cerrado.

\subsection{Bloque A --- Consolidacion del sistema entregado}

Se exige que la totalidad del alcance de la Entrega 2 permanezca operativo, con las siguientes correcciones y refuerzos.

\begin{itemize}
    \item Todos los \emph{endpoints} CRUD de la API REST deben mantenerse funcionando y documentados en la especificacion OpenAPI 3.0 accesible en \texttt{/api/docs}~\cite{fielding2000rest,walls2022spring}.
    \item Toda la autenticacion debe operar bajo cookie \texttt{HttpOnly + Secure + SameSite=Strict} con JWT firmado, siguiendo el flujo definido en el RFC 7519 y con los seis \emph{claims} estandar (\texttt{iss, sub, aud, exp, nbf, iat, jti})~\cite{jones2015jwt}.
    \item Todos los errores deben responder con \emph{ProblemDetails} de acuerdo con el RFC 7807, incluyendo los campos \texttt{type, title, status, detail} y, cuando aplique, \texttt{instance}~\cite{rfc7807}.
    \item El cache Redis debe seguir en el mismo \emph{endpoint} de listado seleccionado en la Entrega 2, ahora con TTL declarado en configuracion externa (no en codigo) y una \emph{hit ratio} objetivo medida empiricamente y reportada en \texttt{docs/mediciones/}.
\end{itemize}

\subsection{Bloque B --- Reproducibilidad automatica desde clonacion limpia}

Este bloque es el mas exigente para el equipo y el mas discriminante para la nota. Su cumplimiento estricto es lo que permite que un revisor externo emita el badge \emph{Artifacts Evaluated --- Reusable} definido por la ACM~\cite{acm2020artifact}.

\begin{cajarepo}{B.1 --- Arranque en un solo comando y sistema autocontenido}
    \begin{itemize}
        \item El repositorio debe contener un \texttt{Makefile} (o \texttt{justfile}, o \texttt{Taskfile.yml}) con al menos los objetivos \texttt{make up}, \texttt{make down}, \texttt{make test}, \texttt{make bench}, \texttt{make audit} y \texttt{make clean}. La sola ejecucion de \texttt{make up} debe levantar el sistema completo sin intervencion humana adicional.
        \item El archivo \texttt{docker-compose.yml} debe pinar cada imagen por \emph{digest} \texttt{sha256} (no solo por \emph{tag}), evitando derivas silenciosas entre reconstrucciones.
        \item Las variables de entorno se centralizan en \texttt{.env.example}, comentadas y con valores por defecto seguros para desarrollo local. El archivo \texttt{.env} real queda excluido por \texttt{.gitignore}.
        \item El esquema de la base de datos se aplica exclusivamente desde \texttt{db/schema.sql} y \texttt{db/seed.sql} montados en \texttt{/docker-entrypoint-initdb.d/}. Queda prohibido depender de \texttt{spring.jpa.hibernate.ddl-auto=update} o similares.
        \item Todo estado observable de la base al inicio se crea con datos semilla reproducibles: el archivo \texttt{db/seed.sql} contiene el usuario \texttt{admin} con hash \emph{BCrypt} valido para la contrasena documentada en el \texttt{README}.
    \end{itemize}
\end{cajarepo}

\begin{cajarepo}{B.2 --- Determinismo de las mediciones}
    \begin{itemize}
        \item Cualquier \emph{script} de generacion de datos, de benchmark o de analisis fija su semilla aleatoria (\texttt{--seed}, \texttt{np.random.seed(42)}, etc.) y la documenta en el propio archivo.
        \item Los benchmarks k6 se ejecutan con la misma configuracion (\texttt{50 VUs}, duracion \texttt{30s}, \emph{ramp-up} declarado en \texttt{k6/opts.js}) para permitir la comparacion entre corridas~\cite{ortega2022securecoding}.
        \item Los reportes Lighthouse se generan con \texttt{npx lhci autorun} contra un contenedor recien levantado, con perfil movil declarado y \emph{throttling} de red \texttt{Slow 4G}.
        \item Cada archivo de resultados incluye en su cabecera la fecha \texttt{ISO 8601}, el \emph{commit hash} corto y las versiones exactas de las herramientas (\texttt{k6 -{-}version}, \texttt{lighthouse -{-}version}, \texttt{java -version}).
    \end{itemize}
\end{cajarepo}

\subsection{Bloque C --- Evidencia empirica cuantitativa}

La calidad de un proyecto de ingenieria de software se mide por evidencia, no por afirmacion. Este bloque exige que cada propiedad no funcional relevante este respaldada por un archivo crudo de mediciones, un \emph{script} de analisis reproducible y un reporte agregado con estadistica descriptiva minima~\cite{wohlin2012experimentation,washizaki2024swebok}.

\subsubsection*{C.1 --- Rendimiento}

Se ejecutan al menos tres corridas independientes de la misma prueba de carga con k6 (\texttt{50 VUs, 30 s}) contra el \emph{endpoint} de listado protegido por la API. Para cada corrida se registran los datos crudos en \texttt{docs/mediciones/perf/kNN-run\{1,2,3\}.json} y se calcula:

\begin{itemize}
    \item Media, mediana, desviacion tipica e intervalo de confianza al $95$\,\% del tiempo de respuesta.
    \item Percentiles \texttt{p50, p90, p95, p99}.
    \item Tasa de errores HTTP $\geq 500$ (debe ser cero).
    \item \emph{Throughput} en peticiones por segundo (media $\pm$ IC 95\,\%).
\end{itemize}

Los umbrales objetivo son: \texttt{p95} de tiempo de respuesta menor a $200\,\text{ms}$ con cache caliente y menor a $500\,\text{ms}$ con cache frio. Estos umbrales estan alineados con los objetivos de calidad recomendados por la ISO/IEC 25010 para la caracteristica de \emph{Eficiencia de desempeno}~\cite{iso25010}.

\subsubsection*{C.2 --- Seguridad}

Se realiza una auditoria automatica de los seis controles OWASP minimos, con evidencia crudo en \texttt{docs/mediciones/sec/}~\cite{owasp2021top10}.

\begin{itemize}
    \item \textbf{A01} (control de acceso roto): \texttt{curl} de un usuario \texttt{A} solicitando el recurso de \texttt{B} debe devolver \texttt{403 Forbidden}. Guardar la salida completa de \texttt{curl -{-}include}.
    \item \textbf{A02} (fallo criptografico): capturar \texttt{curl -v https://localhost:8443} mostrando \texttt{TLSv1.3} y suite AEAD.
    \item \textbf{A03} (inyeccion): enviar \texttt{POST} con payload \texttt{' OR '1'='1} en un campo de busqueda; la respuesta esperada es \texttt{422} con \emph{ProblemDetails}.
    \item \textbf{A05} (mala configuracion de seguridad): \texttt{curl -I} contra la raiz y verificar \texttt{Strict-Transport-Security, X-Frame-Options: DENY, X-Content-Type-Options: nosniff} y \emph{Content-Security-Policy}.
    \item \textbf{A07} (fallo de identificacion y autenticacion): enviar seis intentos fallidos consecutivos de login desde la misma IP y verificar respuesta \texttt{429 Too Many Requests} desde el sexto en adelante.
    \item \textbf{A09} (fallo de registro y monitoreo): capturar el archivo de log donde consten las entradas de \texttt{login} exitoso y fallido con \texttt{ip, timestamp, sub} del JWT.
\end{itemize}

\subsubsection*{C.3 --- Usabilidad}

Se aplica el instrumento \emph{System Usability Scale} de Brooke, con las diez preguntas originales sin modificar la escala de cinco puntos~\cite{brooke1996sus}. El protocolo minimo es:

\begin{itemize}
    \item Poblacion: minimo diez participantes externos al equipo del PFC, con perfil declarado (edad, sexo, experiencia web, dispositivo). El instrumento SUS ha demostrado validez con muestras desde cinco participantes, pero diez es el minimo recomendado para producir estimaciones estables de la media.
    \item Consentimiento informado firmado (fisico o digital) y almacenado en \texttt{docs/etica/consentimientos/} con datos personales anonimizados (los formularios firmados quedan fuera del repositorio publico y se listan por codigo \texttt{P01, P02, \ldots}).
    \item Tarea comun de \emph{onboarding}: login, alta de un registro, edicion, eliminacion logica, logout.
    \item Resultados: matriz cruda (\texttt{docs/mediciones/sus/sus-raw.csv}) con una fila por participante y una columna por item, y reporte agregado con puntuacion media SUS ($0$ a $100$), desviacion tipica e intervalo de confianza al $95$\,\%.
\end{itemize}

\subsubsection*{C.4 --- Cobertura de pruebas}

Se genera reporte JaCoCo con la cobertura de \emph{lines, branches, complexity} sobre el modulo de dominio y de servicios de la API. La cobertura objetivo declarada por el planteamiento del PFC es $\geq 70$\,\% para la Entrega Final; en la Tercera Entrega se exige alcanzar al menos $\geq 60$\,\% y una tendencia creciente entre entregas.

\subsubsection*{C.5 --- Accesibilidad y calidad web}

Se ejecuta \texttt{npx lhci autorun} contra el frontend Angular servido por el contenedor. Los umbrales minimos declarados en \texttt{lighthouserc.js} son:

\begin{itemize}
    \item \emph{Performance} $\geq 80$.
    \item \emph{Accessibility} $\geq 90$.
    \item \emph{Best practices} $\geq 90$.
    \item \emph{SEO} $\geq 90$.
\end{itemize}

El archivo JSON de cada auditoria queda archivado en \texttt{docs/mediciones/lighthouse/} con nombre \texttt{lhci-YYYYMMDD-HHMM.json}.

\subsection{Bloque D --- Documentacion arquitectonica completa y trazable}

La arquitectura documentada en la Entrega 2 se refuerza con los siguientes artefactos.

\begin{itemize}
    \item Diagramas C4 en niveles 1, 2 y 3, con codigo fuente en \emph{Structurizr DSL} versionado en \texttt{docs/arquitectura/} y exportacion PNG generada por el pipeline~\cite{brown2023c4}.
    \item Cinco ADRs siguiendo la plantilla de Nygard (\emph{Title, Context, Decision, Status, Consequences}), archivados en \texttt{docs/adr/adr-NNN-titulo.md}. Los cinco temas obligatorios son: eleccion de la pila principal, esquema de autenticacion, gestor de base de datos, estrategia de cache y estrategia de despliegue~\cite{nygard2011adr}.
    \item Tabla de atributos de calidad ISO/IEC 25010, con prioridad, escenario y estrategia por cada uno~\cite{iso25010}.
    \item Matriz de trazabilidad \emph{requisitos $\rightarrow$ modulos $\rightarrow$ endpoints $\rightarrow$ pruebas $\rightarrow$ evidencia} en \texttt{docs/trazabilidad/matriz.csv} o \texttt{.md}.
\end{itemize}

\subsection{Bloque E --- Publicabilidad del artefacto y gestion de datos}

Este bloque orienta la Tercera Entrega hacia el estado de reutilizabilidad exigido por las guias FAIR y por el badge \emph{Artifacts Available} de la ACM~\cite{wilkinson2016fair,acm2020artifact}. Su cumplimiento se demuestra con archivos concretos dentro del repositorio.

\begin{cajaentregable}{E.1 --- Licencia y metadatos de citacion}
    \begin{itemize}
        \item \texttt{LICENSE} en la raiz con el texto integro de una licencia aprobada por la Open Source Initiative (MIT, Apache-2.0 o BSD-3-Clause). La eleccion se justifica en un ADR.
        \item \texttt{CITATION.cff} en la raiz siguiendo la version \texttt{1.2.0} de la especificacion, con los campos obligatorios (\texttt{cff-version, message, title, authors}) y los recomendados (\texttt{version, date-released, license, repository-code, doi, keywords})~\cite{druskat2021cff,smith2016software}.
        \item Las contribuciones de cada integrante se declaran en \texttt{CONTRIBUTORS.md} usando la taxonomia CRediT (14 roles), asignando explicitamente a cada persona los roles en que participo~\cite{niso2022credit}.
    \end{itemize}
\end{cajaentregable}

\begin{cajaentregable}{E.2 --- Identificador persistente y archivo del artefacto}
    \begin{itemize}
        \item Se crea el archivo Zenodo del repositorio en el estado \texttt{v0.9.0-rc}, obteniendo un DOI persistente. El DOI se documenta en el \texttt{README}, en \texttt{CITATION.cff} y en el propio archivo del PDF de informe.
        \item El badge Zenodo (\texttt{[![DOI](\ldots)]}) aparece al inicio del \texttt{README}, junto al badge de \emph{CI passing} y al badge de licencia.
    \end{itemize}
\end{cajaentregable}

\begin{cajaentregable}{E.3 --- Diccionario de datos}
    \begin{itemize}
        \item Para cada archivo crudo de mediciones bajo \texttt{docs/mediciones/} se registra en \texttt{docs/mediciones/DATA-DICTIONARY.md} el nombre de la variable, su tipo de dato, la unidad, el rango esperado y su significado.
        \item El diccionario se actualiza con \emph{pull request} cada vez que se anade una nueva variable, y su vigencia se verifica automaticamente en el pipeline con un test de esquema (por ejemplo con \emph{jsonschema} o \emph{pandera} si se usa Python).
    \end{itemize}
\end{cajaentregable}

\begin{cajaentregable}{E.4 --- Versionado semantico y \emph{changelog}}
    \begin{itemize}
        \item El proyecto adopta \emph{Semantic Versioning 2.0.0} de forma estricta desde la Tercera Entrega. La regla \texttt{MAJOR.MINOR.PATCH} se documenta en \texttt{docs/VERSIONING.md} y el tag de esta entrega debe ser exactamente \texttt{v0.9.0-rc}~\cite{prestonwerner2013semver}.
        \item El archivo \texttt{CHANGELOG.md} sigue la convencion \emph{Keep a Changelog}, con secciones \texttt{Added, Changed, Deprecated, Removed, Fixed, Security} para cada version.
        \item Los mensajes de \emph{commit} siguen \emph{Conventional Commits} (\texttt{feat:, fix:, docs:, chore:, refactor:, test:, perf:}) para permitir la generacion automatica del changelog.
    \end{itemize}
\end{cajaentregable}

\subsection{Bloque F --- Etica y responsabilidad de los datos}

Aunque el proyecto no manipule datos de salud, financieros o personales sensibles, la disciplina de investigacion contemporanea exige declarar explicitamente los aspectos eticos del uso de datos y de participantes~\cite{washizaki2024swebok,wohlin2012experimentation}.

\begin{itemize}
    \item \texttt{docs/etica/ETHICS.md} declara: (i) fuentes de datos y su licencia, (ii) tratamiento de datos personales si los hubiera, (iii) mecanismo de consentimiento informado para las pruebas de usabilidad, y (iv) ausencia de datos identificables en el repositorio publico.
    \item \texttt{docs/etica/consentimientos/plantilla.md} contiene la plantilla de consentimiento informado utilizada en las pruebas SUS del Bloque C.3. Los consentimientos firmados no se suben al repositorio publico.
\end{itemize}

% =====================================================================
%  3. ESTRUCTURA DE REPOSITORIO
% =====================================================================
\section{Estructura obligatoria del repositorio para la Tercera Entrega}

El arbol de directorios que el revisor debe encontrar al clonar el repositorio en el estado \texttt{v0.9.0-rc} es el siguiente. Cualquier desviacion se considera evidencia de una entrega incompleta y se penaliza en el criterio de calidad de entrega.

\begin{lstlisting}[caption={Estructura minima del repositorio para la Tercera Entrega}]
.
|-- README.md                   # badges, arranque, credenciales,
|                               # DOI Zenodo, licencia
|-- LICENSE                     # texto integro OSI-approved
|-- CITATION.cff                # metadatos de citacion CFF 1.2.0
|-- CONTRIBUTORS.md             # roles CRediT por integrante
|-- CHANGELOG.md                # Keep-a-Changelog
|-- Makefile                    # objetivos up down test bench audit
|-- docker-compose.yml          # con digests sha256 pinados
|-- .env.example                # variables comentadas
|-- .gitignore
|-- backend/                    # codigo Spring Boot
|-- frontend/                   # codigo Angular
|-- db/
|   |-- schema.sql
|   `-- seed.sql
|-- docs/
|   |-- adr/                    # 5 ADRs plantilla Nygard
|   |   |-- adr-001-...md
|   |   |-- ...
|   |   `-- adr-005-...md
|   |-- arquitectura/           # Structurizr DSL + PNG C4 L1..L3
|   |-- mediciones/
|   |   |-- perf/               # k6 raw runs 1..3 + reporte
|   |   |-- sec/                # curl -I, curl -v, capturas OWASP
|   |   |-- sus/                # SUS raw csv + reporte + IC 95%
|   |   |-- lighthouse/         # JSON lhci por corrida
|   |   |-- jacoco/             # report.xml + report.html
|   |   `-- DATA-DICTIONARY.md  # variables + tipos + unidades
|   |-- trazabilidad/
|   |   `-- matriz.csv          # req -> modulo -> endpoint -> test
|   |                           #    -> evidencia
|   |-- etica/
|   |   |-- ETHICS.md
|   |   `-- consentimientos/plantilla.md
|   |-- VERSIONING.md
|   `-- informe-entrega-3.pdf   # informe tecnico 15-25 paginas
|-- k6/                         # scripts + opts.js
|-- lighthouserc.js
|-- scripts/                    # analisis con seed fija
`-- .github/workflows/          # CI que ejecuta test bench audit
\end{lstlisting}

% =====================================================================
%  4. ENTREGABLES
% =====================================================================
\section{Entregables consolidados de la Tercera Entrega}

El equipo entrega, unicamente a traves del repositorio Git publico y del formulario institucional del SGA, los siguientes artefactos.

\begin{enumerate}
    \item Repositorio Git publico con \emph{tag} exactamente igual a \texttt{v0.9.0-rc}, cuyo \emph{commit hash} corto se declara en la portada del PDF de informe tecnico.
    \item DOI persistente asignado por Zenodo al tag \texttt{v0.9.0-rc}. El DOI se declara en tres lugares (README, CITATION.cff, portada del PDF) y debe resolver a la misma version.
    \item PDF de informe tecnico de 20 a 30 paginas, ubicado en \texttt{docs/informe-entrega-3.pdf}, con las secciones: portada, resumen ejecutivo, estado del sistema, arquitectura actual (C4 L1--L3), matriz de trazabilidad, protocolo experimental, resultados por bloque (rendimiento, seguridad, usabilidad, cobertura, accesibilidad), amenazas a la validez, declaracion de contribuciones (CRediT), declaracion etica, referencias.
    \item Coleccion Postman actualizada (\texttt{docs/postman/coleccion.json}) con al menos 20 peticiones cubriendo casos de exito, error de validacion (\texttt{422}), no autorizado (\texttt{401}), prohibido (\texttt{403}) y no encontrado (\texttt{404}).
    \item Video corto (2--3 minutos, alojado enlazado desde el README) donde el equipo demuestra \texttt{make up}, ejecucion de \texttt{make bench} y \texttt{make audit}, y muestra los reportes generados. No sustituye a la evidencia archivada.
\end{enumerate}

% =====================================================================
%  5. RUBRICA
% =====================================================================
\section{Rubrica de evaluacion de la Tercera Entrega}

La rubrica pondera diez criterios agrupados en dos ejes: eje de ingenieria (criterios C1--C6, $60$\,\%) y eje de publicabilidad del artefacto (criterios C7--C10, $40$\,\%). Cada criterio se evalua en cinco niveles: \emph{Excelente} $= 100$\,\%, \emph{Satisfactorio} $= 75$\,\%, \emph{En desarrollo} $= 50$\,\%, \emph{Insuficiente} $= 25$\,\%, \emph{Ausente} $= 0$\,\%. La escala de cuatro niveles activos mas Ausente ya fue empleada por el docente en la rubrica GA-PE-U2 y se mantiene aqui por consistencia interna del curso.

\subsection{Formula de calificacion}

\begin{center}
$\displaystyle \text{Nota}_{100} \;=\; \sum_{i=1}^{10} \bigl(\text{peso}_{i} \times \text{nivel}_{i}\bigr) \qquad \text{Nota}_{10} \;=\; \frac{\text{Nota}_{100}}{10}$
\end{center}

\subsection{Rubrica detallada}

{
\footnotesize
\renewcommand{\arraystretch}{1.3}
\begin{longtable}{|L{3.0cm}|C{0.9cm}|L{2.5cm}|L{2.5cm}|L{2.5cm}|L{2.3cm}|}
\hline
\rowcolor{uteqverde!85}
{\color{white}\textbf{Criterio}} &
{\color{white}\textbf{Peso}} &
{\color{white}\textbf{Excelente (100\,\%)}} &
{\color{white}\textbf{Satisfactorio (75\,\%)}} &
{\color{white}\textbf{En desarrollo (50\,\%)}} &
{\color{white}\textbf{Insuficiente (25\,\%)}} \\
\hline
\endfirsthead

\hline
\rowcolor{uteqverde!85}
{\color{white}\textbf{Criterio}} &
{\color{white}\textbf{Peso}} &
{\color{white}\textbf{Excelente (100\,\%)}} &
{\color{white}\textbf{Satisfactorio (75\,\%)}} &
{\color{white}\textbf{En desarrollo (50\,\%)}} &
{\color{white}\textbf{Insuficiente (25\,\%)}} \\
\hline
\endhead

% ============ C1 ============
\textbf{C1. Consolidacion funcional} (Bloque A) &
\centering 8\,\% &
Todos los \emph{endpoints} de la Entrega 2 operativos, JWT en cookie HttpOnly+Secure+SameSite, ProblemDetails RFC 7807 en el 100\,\% de errores, cache Redis operativo con TTL externo. &
Igual pero con una discrepancia menor documentada. &
Al menos un \emph{endpoint} caido o ProblemDetails en menos del 80\,\% de errores. &
Regresiones funcionales significativas respecto a la Entrega 2. \\
\rowcolor{grisfila}

% ============ C2 ============
\textbf{C2. Reproducibilidad automatica} (Bloque B.1) &
\centering 10\,\% &
\texttt{make up} arranca todo desde clonacion limpia sin intervencion; imagenes pinadas por digest sha256; \texttt{.env.example} completo; seed y schema montados en \texttt{docker-entrypoint-initdb.d}. &
Arranque en un solo comando con una variable de entorno adicional documentada. &
Requiere pasos manuales no documentados o falla ocasional. &
No arranca desde clonacion limpia siguiendo el README. \\

% ============ C3 ============
\textbf{C3. Determinismo de mediciones} (Bloque B.2) &
\centering 7\,\% &
Todos los scripts con semilla fija y versiones de herramientas capturadas en cada archivo de resultados; 3 corridas independientes por benchmark. &
Semillas fijas en scripts principales; menos de 3 corridas o versiones parcialmente registradas. &
Semillas presentes en algunos scripts; corridas unicas. &
Sin semillas fijas ni control de versiones de herramientas. \\
\rowcolor{grisfila}

% ============ C4 ============
\textbf{C4. Evidencia empirica cuantitativa} (Bloque C) &
\centering 12\,\% &
Los cinco sub-bloques (perf, sec, sus, cobertura, lighthouse) con archivos crudos, scripts de analisis y reportes con media, IC 95\,\% y percentiles; SUS con $\geq 10$ participantes. &
Cuatro de los cinco sub-bloques completos; el faltante justificado en el informe. &
Tres de los cinco sub-bloques; SUS con menos de 10 participantes. &
Menos de tres sub-bloques con evidencia; SUS ausente. \\

% ============ C5 ============
\textbf{C5. Documentacion arquitectonica y trazabilidad} (Bloque D) &
\centering 10\,\% &
C4 L1--L3 versionado en Structurizr DSL; cinco ADRs con plantilla Nygard completa; tabla ISO 25010; matriz de trazabilidad completa \emph{req $\rightarrow$ codigo $\rightarrow$ prueba $\rightarrow$ evidencia}. &
Cuatro ADRs completos; matriz cubre $\geq 80$\,\% de requisitos. &
Tres ADRs; matriz parcial. &
Menos de tres ADRs y matriz ausente o superficial. \\
\rowcolor{grisfila}

% ============ C6 ============
\textbf{C6. Auditoria de seguridad OWASP} (Bloque C.2 con enfasis) &
\centering 13\,\% &
Los seis controles OWASP evidenciados con \texttt{curl} verificable, respuestas HTTP correctas (403/422/429), cabeceras HSTS, CSP, X-Frame-Options, X-Content-Type-Options presentes; logs con timestamp, IP y sub. &
Cinco controles evidenciados; una cabecera menor faltante. &
Cuatro controles evidenciados. &
Tres o menos controles evidenciados. \\

% ============ C7 ============
\textbf{C7. Licencia, citacion y contribucion} (Bloque E.1) &
\centering 8\,\% &
\texttt{LICENSE} OSI-approved presente; \texttt{CITATION.cff} v1.2.0 valido; \texttt{CONTRIBUTORS.md} con roles CRediT explicitos por integrante. &
Licencia y CITATION.cff presentes; roles CRediT parciales. &
Licencia presente; CITATION.cff ausente o invalido. &
Sin licencia declarada. \\
\rowcolor{grisfila}

% ============ C8 ============
\textbf{C8. Archivo permanente e identificador} (Bloque E.2) &
\centering 7\,\% &
DOI Zenodo asignado al tag \texttt{v0.9.0-rc}, declarado en tres lugares y resolviendo a la misma version. &
DOI asignado y declarado en dos lugares. &
DOI asignado pero no declarado en el repositorio. &
Sin DOI persistente. \\

% ============ C9 ============
\textbf{C9. Datos, diccionario y versionado} (Bloques E.3, E.4) &
\centering 10\,\% &
Datos crudos en \texttt{docs/mediciones/}; \texttt{DATA-DICTIONARY.md} completo; SemVer 2.0.0 estricto; \texttt{CHANGELOG.md} con Keep-a-Changelog; commits siguen Conventional Commits. &
Diccionario cubre $\geq 80$\,\% de variables; changelog presente. &
Diccionario parcial; changelog irregular. &
Sin diccionario o sin changelog. \\
\rowcolor{grisfila}

% ============ C10 ============
\textbf{C10. Etica y calidad de la entrega global} (Bloque F + integridad global) &
\centering 15\,\% &
\texttt{ETHICS.md} y plantilla de consentimiento presentes; PDF de informe con las 10 secciones exigidas; matriz de amenazas a la validez presente; historial Git con $\geq 30$ commits granulares y autoria correcta de todos los integrantes. &
Documentacion etica presente; informe con una seccion faltante justificada. &
\texttt{ETHICS.md} minimo; informe con dos secciones faltantes. &
Sin declaracion etica ni informe estructurado. \\

\hline
\rowcolor{goldclaro}
\multicolumn{2}{|r|}{\textbf{Total ponderado}} &
\multicolumn{4}{l|}{\textbf{100\,\%} \quad \footnotesize Nota${}_{100} = \sum (\text{peso}_i \times \text{nivel}_i)$ \; \textbar\; Nota${}_{10} = \text{Nota}_{100}/10$} \\
\hline
\end{longtable}
}

\subsection{Reglas transversales de calificacion}

\begin{cajacritica}{Reglas transversales que aplican a todos los criterios}
    \begin{enumerate}
        \item \textbf{El repositorio Git es la fuente de verdad}. La evidencia verificada directamente en el repositorio se distingue siempre de la inferida de un PDF o de capturas.
        \item Los criterios dependientes de codigo (C1--C6) se califican \emph{Ausente} ($0$\,\%) cuando no existe repositorio publico accesible o cuando el tag \texttt{v0.9.0-rc} no existe.
        \item Cualquier commit posterior al corte declarado en la primera pagina del informe se ignora para efectos de calificacion, pero no se penaliza (queda como evidencia de trabajo posterior a la entrega).
        \item Si el \texttt{make up} falla desde una clonacion limpia siguiendo estrictamente el README, el criterio C2 se califica automaticamente \emph{Insuficiente} ($25$\,\%) o menos, independientemente del resto.
        \item Si el DOI declarado no resuelve al tag \texttt{v0.9.0-rc}, el criterio C8 se califica \emph{Ausente} ($0$\,\%).
        \item Si el \texttt{CITATION.cff} no valida contra el esquema \texttt{cff-version: 1.2.0} publicado en Zenodo (validado con \texttt{cffconvert}), el criterio C7 se limita a \emph{En desarrollo} ($50$\,\%) como maximo~\cite{druskat2021cff}.
        \item Todos los archivos crudos de mediciones deben conservarse tal cual los produjo la herramienta original; su edicion manual invalida la evidencia.
    \end{enumerate}
\end{cajacritica}

% =====================================================================
%  6. TRABAJO PENDIENTE PARA LA CUARTA ENTREGA
% =====================================================================
\section{Consideraciones para la Entrega Final (semana 16)}

La Tercera Entrega deja al equipo en el estado \texttt{v0.9.0-rc}, es decir, un candidato a version estable. La transicion a \texttt{v1.0.0} en la Entrega Final consistira en absorber las observaciones de la evaluacion de esta entrega, elevar la cobertura JaCoCo desde el minimo del $60$\,\% hasta el objetivo del $70$\,\% o superior, completar las mediciones que la Tercera Entrega haya dejado incompletas, ampliar la muestra SUS si aplica y consolidar el archivo permanente. En consecuencia, ninguna decision arquitectonica mayor debe posponerse a la Entrega Final: si algo requiere reescribir un ADR o refactorizar un modulo, debe ocurrir dentro de la Tercera Entrega.

% =====================================================================
%  REFERENCIAS
% =====================================================================
\bibliographystyle{plain}
\bibliography{Entrega3}

\end{document}
